\chapter{Chapter09 — Improved KKL}
%%% generated by the blueprint pipeline from FABL.Chapter09.ImprovedKKL
\section{Overview}

Declarations of the Lean module \texttt{FABL.Chapter09.ImprovedKKL}.

\section{Declarations}

\begin{theorem}[Variance and total stable influence]
\lean{FABL.variance_mul_rho_rpow_sub_one_le_totalStableInfluence}
\proofsource{Analysis-of-Boolean-Functions-by-Ryan-ODonnell}{
  pdf-ffbe9c7b3306-p267-b005,
  pdf-ffbe9c7b3306-p268-b001
}
\label{FABL.variance_mul_rho_rpow_sub_one_le_totalStableInfluence}
\leanok
\uses{FABL.BooleanFunction, FABL.BooleanFunction.toReal, FABL.fourierCoeff, FABL.totalInfluence, FABL.totalInfluence_eq_sum_card_mul_sq_fourierCoeff, FABL.totalStableInfluence, FABL.totalStableInfluence_eq_sum_card_mul_rho_pow_mul_sq_fourierCoeff, FABL.variance, FABL.variance_eq_sum_sq_fourierCoeff}
\difficulty{6}
Let $n$ be a natural number and let $f:\{−1,1\}^n\to\{−1,1\}$ be a Boolean function with
positive variance. For every $\rho\in(0,1]$, its variance multiplied by
$\rho^{I(f)/\operatorname{Var}(f)-1}$ is at most its total stable influence at parameter
$\rho$, where $I(f)$ denotes its total influence.
\end{theorem}

\begin{theorem}[Stable influence upper bound]
\lean{FABL.totalStableInfluence_le_maximumInfluence_rpow_mul}
\label{FABL.totalStableInfluence_le_maximumInfluence_rpow_mul}
\leanok
\uses{FABL.BooleanFunction, FABL.BooleanFunction.toReal, FABL.booleanInfluence, FABL.booleanInfluence_eq_influence_toReal, FABL.booleanInfluence_le_maximumInfluence, FABL.influence, FABL.influence_nonneg, FABL.maximumInfluence, FABL.maximumInfluence_nonneg, FABL.stableInfluence_le_booleanInfluence_rpow, FABL.totalInfluence, FABL.totalStableInfluence}
\difficulty{4}
Let $n$ be a nonnegative integer and let $f:\{-1,1\}^n\to\{-1,1\}$ be a Boolean
function. If $\delta\ge 0$, $0\le\rho\le 1$, and $2/(1+\rho)=1+\delta$, then the
total stable influence of the real-valued sign representation of $f$ at parameter
$\rho$ is at most
$\operatorname{Inf}_{\max}(f)^\delta\,\operatorname{Inf}_{\mathrm{tot}}(f)$, where
$\operatorname{Inf}_{\max}$ is the maximum coordinate influence and
$\operatorname{Inf}_{\mathrm{tot}}$ is the total influence of that sign
representation.
\end{theorem}

\begin{definition}[Exercise 9.30 correlation parameter]
\lean{FABL.exercise9_30Rho}
\label{FABL.exercise9_30Rho}
\leanok
For a real number $\delta$, define $\rho=(1-\delta)/(1+\delta)$.
\end{definition}

\begin{theorem}[Range of the ratio parameter]
\lean{FABL.exercise9_30Rho_mem_Ioc}
\label{FABL.exercise9_30Rho_mem_Ioc}
\leanok
\uses{FABL.exercise9_30Rho}
\difficulty{3}
For every real number $\delta$ with $0<\delta<1$, the number
$\rho=(1-\delta)/(1+\delta)$ satisfies $0<\rho\le1$.
\end{theorem}

\begin{theorem}[Exponent identity for the ratio parameter]
\lean{FABL.exercise9_30_exponent_identity}
\label{FABL.exercise9_30_exponent_identity}
\leanok
\uses{FABL.exercise9_30Rho}
\difficulty{3}
For every real number $\delta$ with $0<\delta<1$, if $\rho=(1-\delta)/(1+\delta)$, then
$2/(1+\rho)=1+\delta$.
\end{theorem}

\begin{theorem}[Compact maximum-influence bound]
\lean{FABL.exercise9_30_compact}
\label{FABL.exercise9_30_compact}
\leanok
\uses{FABL.BooleanFunction, FABL.BooleanFunction.toReal, FABL.edgeKKL, FABL.exercise9_30Rho, FABL.exercise9_30Rho_mem_Ioc, FABL.exercise9_30_exponent_identity, FABL.maximumInfluence, FABL.totalInfluence, FABL.totalStableInfluence_le_maximumInfluence_rpow_mul, FABL.variance, FABL.variance_le_totalInfluence, FABL.variance_mul_rho_rpow_sub_one_le_totalStableInfluence}
\difficulty{5}
Let $f:\{−1,1\}^n\to\{−1,1\}$ have positive variance, and let $0<\delta<1$. Write $I$
for its total influence, $V$ for its variance, and $\rho=(1-\delta)/(1+\delta)$. Then
$\left(\rho^{I/V-1}/(I/V)\right)^{1/\delta}\le \max_i\operatorname{Inf}_i(f)$, with
variance and coordinate influences computed under the uniform sign-cube distribution.
\end{theorem}

\begin{theorem}[Lower bound for maximum influence]
\lean{FABL.exercise9_30_rho_form}
\label{FABL.exercise9_30_rho_form}
\leanok
\uses{FABL.BooleanFunction, FABL.BooleanFunction.toReal, FABL.exercise9_30Rho, FABL.exercise9_30Rho_mem_Ioc, FABL.exercise9_30_compact, FABL.maximumInfluence, FABL.totalInfluence, FABL.variance, FABL.variance_le_totalInfluence}
\difficulty{4}
Let $f$ be a Boolean function on the discrete cube, and suppose its variance is
positive. For $0<\delta<1$, define $\rho=(1-\delta)/(1+\delta)$. Then
\[
\rho^{-1/\delta}\left(\frac{I(f)}{\operatorname{Var}(f)}\right)^{-1/\delta}\rho^{I(f)/(\delta\operatorname{Var}(f))}\leq
\max_i \operatorname{Inf}_i(f),
\]
where $I(f)$ is the total influence and $\operatorname{Inf}_i(f)$ is the influence of
coordinate $i$.
\end{theorem}

\begin{theorem}[Parameterized maximum-influence bound]
\lean{FABL.exercise9_30_parameterized}
\label{FABL.exercise9_30_parameterized}
\leanok
\uses{FABL.BooleanFunction, FABL.BooleanFunction.toReal, FABL.exercise9_30Rho, FABL.exercise9_30_rho_form, FABL.maximumInfluence, FABL.totalInfluence, FABL.variance}
\difficulty{3}
Let $f:\{−1,1\}^n\to\{−1,1\}$ have positive variance, and let $0<\delta<1$. Put $K=I/V$,
where $I$ is the total influence of $f$ and $V$ its variance, both for the uniform
sign-cube distribution. Then
$\left(\frac{1+\delta}{1-\delta}\right)^{1/\delta}\left(\frac1K\right)^{1/\delta}\left(\frac{1-\delta}{1+\delta}\right)^{K/\delta}\le\max_i\operatorname{Inf}_i(f)$.
\end{theorem}

\begin{definition}[Exercise 9.30 logarithmic remainder]
\lean{FABL.exercise9_30LogRemainder}
\label{FABL.exercise9_30LogRemainder}
\leanok
For a real number $\delta$, define the logarithmic remainder by
$\log(1-\delta)-\log(1+\delta)+2\delta+\delta^3$.
\end{definition}

\begin{theorem}[Nonnegative logarithmic remainder]
\lean{FABL.exercise9_30LogRemainder_nonneg}
\label{FABL.exercise9_30LogRemainder_nonneg}
\leanok
\uses{FABL.exercise9_30LogRemainder}
\difficulty{6}
For every real number $\delta$ with $0\le\delta\le\tfrac12$,
$\log(1-\delta)-\log(1+\delta)+2\delta+\delta^3\ge 0$.
\end{theorem}

\begin{theorem}[Lower bound for a power of rho]
\lean{FABL.exercise9_30_rho_rpow_inv_lower}
\label{FABL.exercise9_30_rho_rpow_inv_lower}
\leanok
\uses{FABL.exercise9_30LogRemainder, FABL.exercise9_30LogRemainder_nonneg, FABL.exercise9_30Rho}
\difficulty{4}
For $0<\delta\leq 1/2$, let $\rho=(1-\delta)/(1+\delta)$. Then
\[
\exp(-2-\delta^2)\leq \rho^{1/\delta}.
\]
\end{theorem}

\begin{theorem}[Refined maximum-influence bound]
\lean{FABL.exercise9_30_refined}
\label{FABL.exercise9_30_refined}
\leanok
\uses{FABL.BooleanFunction, FABL.BooleanFunction.toReal, FABL.exercise9_30Rho, FABL.exercise9_30Rho_mem_Ioc, FABL.exercise9_30_compact, FABL.exercise9_30_rho_rpow_inv_lower, FABL.maximumInfluence, FABL.totalInfluence, FABL.variance, FABL.variance_le_totalInfluence}
\difficulty{5}
Let $f:\{−1,1\}^n\to\{−1,1\}$ have positive variance, and put $K=I/V$, where $I$ is its
total influence and $V$ its variance under the uniform sign-cube distribution. Then
$\exp(-2K)\exp(2)K^{-1}{}^{2K^{1/3}}\exp(-\tfrac14K^{1/3})\le\max_i\operatorname{Inf}_i(f)$.
\end{theorem}

\begin{definition}[Logarithmic error scale]
\lean{FABL.exercise9_30KKLError}
\statementsource{Analysis-of-Boolean-Functions-by-Ryan-ODonnell}{
  pdf-ffbe9c7b3306-p267-b003,
  pdf-ffbe9c7b3306-p267-b005,
  pdf-ffbe9c7b3306-p268-b001
}
\label{FABL.exercise9_30KKLError}
\leanok
For each natural number $n$, define the error scale by $1/\sqrt{\log n}$, where $n$ is
viewed as a real number inside the logarithm.
\end{definition}

\begin{theorem}[Asymptotic lower bound on maximum influence]
\lean{FABL.exercise9_30_asymptotic_bound_of_large}
\label{FABL.exercise9_30_asymptotic_bound_of_large}
\leanok
\uses{FABL.BooleanFunction, FABL.BooleanFunction.toReal, FABL.exercise9_30KKLError, FABL.exercise9_30_refined, FABL.maximumInfluence, FABL.maximumInfluence_nonneg, FABL.totalInfluence, FABL.totalInfluence_le_natCast_mul_maximumInfluence, FABL.variance, FABL.variance_eq_four_mul_probabilities, FABL.variance_le_totalInfluence}
\difficulty{6}
Let $n$ be a natural number such that $1<\log n$, $96(\log n)^{5/12}\le\sqrt{\log n}$,
and $(\log n)/2\le\exp(\sqrt{\log n}/2)$. For every function $f:\{−1,1\}^n\to\{−1,1\}$,
its maximum coordinate influence is at least $\frac12\operatorname{Var}(f)\frac{\log
n}{n}\bigl(1-(\log n)^{-1/2}\bigr)$, where the variance and influences are taken with
respect to the uniform sign cube.
\end{theorem}

\begin{theorem}[Vanishing reciprocal square root of the logarithm]
\lean{FABL.tendsto_exercise9_30KKLError}
\label{FABL.tendsto_exercise9_30KKLError}
\leanok
\uses{FABL.exercise9_30KKLError}
\difficulty{3}
For each natural number $n$, let $a_n = 1/\sqrt{\log(n : \mathbb{R})}$, where the
reciprocal is taken in $\mathbb{R}$. Then $a_n$ tends to $0$ as $n$ tends to infinity.
\end{theorem}

\begin{theorem}[Eventual numerical conditions]
\lean{FABL.eventually_exercise9_30_numeric_conditions}
\label{FABL.eventually_exercise9_30_numeric_conditions}
\leanok
\difficulty{6}
For all sufficiently large natural numbers $n$, the following inequalities hold: $1<\log
n$, $96(\log n)^{5/12}\le\sqrt{\log n}$, and $(\log n)/2\le\exp(\sqrt{\log n}/2)$.
\end{theorem}

\begin{theorem}[Eventually exercise9 30 asymptotic kkl]
\lean{FABL.eventually_exercise9_30_asymptotic_kkl}
\statementsource{Analysis-of-Boolean-Functions-by-Ryan-ODonnell}{
  pdf-ffbe9c7b3306-p267-b003,
  pdf-ffbe9c7b3306-p267-b005,
  pdf-ffbe9c7b3306-p268-b001
}
\label{FABL.eventually_exercise9_30_asymptotic_kkl}
\leanok
\uses{FABL.BooleanFunction, FABL.BooleanFunction.toReal, FABL.eventually_exercise9_30_numeric_conditions, FABL.exercise9_30KKLError, FABL.exercise9_30_asymptotic_bound_of_large, FABL.maximumInfluence, FABL.variance}
\difficulty{2}
% dataset-pipeline: stub -- prose pending build_informal_statements.py
Statement pending regeneration.
\end{theorem}
